参数扰动代价与PAC--Bayes中的两项控制。上部示意窄谷、宽谷附近的参数扰动与随机预测器；在相应定理条件下，泛化保证还需同时控制后验经验风险与相对于独立先验的相对熵代价。下部两个面板共用相同横纵尺度，横轴为同一单位方向上的参数位移$s$，即$\bm\theta=\hat{\bm\theta}+s\bm u$、$\|\bm u\|_2=1$。构造的窄谷与宽谷分别为$0.1+0.45s^2$、$0.1+0.08s^2$；实心点为谷底，空心点对应相同位移$s=0.6$，经验风险分别增加0.162和0.0288。上部轮廓没有数值含义，下部曲线仅为可复算的教学构造。谷的宽度受参数化与坐标尺度影响，宽谷不必具有较小的期望风险。
